% =======> To use this template:
% => Edit authors in this file.
% => Edit the paper's content itself in the respective sections. Edit as needed
% => Add any references to references.bib (cite with \cite or \citet)
% =================================================

\documentclass{article}

% to avoid loading the natbib package, add option nonatbib:
% \usepackage[nonatbib]{neurips_2024}
% \PassOptionsToPackage{square,comma,sort,numbers}{natbib}
\PassOptionsToPackage{numbers, compress}{natbib}
% \usepackage[numbers]{natbib}

% to compile a preprint version, e.g., for submission to arXiv, add add the
% [preprint] option:
\usepackage[preprint]{neurips_2024}
% \usepackage[nonatbib, preprint]{neurips_2024}

\usepackage[utf8]{inputenc} % allow utf-8 input
% \usepackage[T1]{fontenc}    % use 8-bit T1 fonts
% For maths/theorems and such
\usepackage{amsmath}
\usepackage{amssymb}
\usepackage{mathtools}
\usepackage{amsthm}
\usepackage{float}
% Use DeclareMathOperator and/or \newcommand to define abbreviations, e.g.:
\DeclareMathOperator{\R}{\mathbb{R}} % real numbers
% Recommended, but optional, packages for figures and better typesetting:
\usepackage{microtype}
\usepackage{graphicx}
\usepackage{subfigure}
\usepackage{subcaption}
\usepackage{booktabs} % for professional tables (https://www.ctan.org/pkg/booktabs)
\usepackage{algorithm}  % For algorithms
\usepackage{algorithmic}
\usepackage{hyperref} % hyperref makes hyperlinks in the resulting PDF.
\usepackage{wrapfig}  % For wrapping figures with text
\usepackage[capitalize,noabbrev]{cleveref}   % For \cref
\RequirePackage{doi}
% Color references and URLs
\usepackage{xcolor, colortbl}
\definecolor{darkblue}{rgb}{0.19, 0.19, 0.62}
\hypersetup{
  colorlinks = true,
  citecolor  = darkblue,
  filecolor  = darkblue,
  urlcolor   = darkblue,
    % linkcolor  = myred,
}
% Todonotes is useful during development; simply uncomment the next line
%    and comment out the line below the next line to turn off comments
%\usepackage[disable,textsize=tiny]{todonotes}
\usepackage[textsize=tiny]{todonotes}

% \input{_macros}



\title{
CSE 151B/251B Milestone Report
\\
% {\small $<$Your-accessible-and-up-to-date-GitHub-URL$>$} % Note: You can put it in your abstracts or as a footnote in the first page, if you prefer. Note 2: Use \href or \url to link your URL.
} 

\author{%
  Haris Zia
  % \thanks{Use footnote for providing further information about author (webpage, alternative address)}
  \\
  Department of Computer Science\\
  University of California San Diego\\
  \texttt{harzia@ucsd.edu} \\
  % examples of more authors
  \And
  Blake Newhouse \\
  Department of Computer Science\\
  University of California San Diego\\
  \texttt{blnewhouse@ucsd.edu} \\
  \AND
  Srideep Dornala \\
  Department of Computer Science\\
  University of California San Diego\\
  \texttt{sdornala@ucsd.edu} \\
  \And
  Daniel Ahossou\\
  Department of Computer Science\\
  University of California San Diego\\
  \texttt{dahossou@ucsd.edu} \\
}
\begin{document}
\maketitle


\begin{abstract}
\textcolor{red}{TODO} Check out \href{https://www.alignmentforum.org/posts/eJGptPbbFPZGLpjsp/highly-opinionated-advice-on-how-to-write-ml-papers#Abstract}{this blog post for tips on how to write a compelling abstract and paper.} \textbf{\textit{Include your final private Kaggle leaderboard score and/or rank at the end of the abstract or introduction.}}
\end{abstract}

\section*{Formatting instructions (remove this section in your submission)}
Please read the instructions below carefully and follow them faithfully. You can keep the section headers but delete the instructional texts.
\paragraph{Hints on writing}
\begin{itemize}
    \item Each paragraph should convey one main message, typically in the beginning sentence of the paragraph (topical sentence).
    \item  Use simple sentences with fewer compounds or clauses. The simpler, the better. Avoid passive tense and use active tense.
    \item Use PDF for your figures to ensure the resolution of your images.
    \item The captions of Figures and Tables should contain sufficient details to be self-explanatory (readers can understand the figure without referring to the text).
\end{itemize}


Figures and tables must be appropriately referenced and discussed in the main text (tip: Use \textbackslash label\{fig:xyz\} and \textbackslash cref\{fig:xyz\}; example:~\cref{fig:example}). Captions must strive to be self-contained.
\paragraph{Figures.}
% \begin{figure}[h]  % Use figure or wrapfigure as needed.
\begin{wrapfigure}{r}{0.4\textwidth} % Example: occupy 40% of text width
  \vspace{-4mm}  % Force moving the figure up by a little bit
  \centering
  \fbox{\rule[-.5cm]{0cm}{4cm} \rule[-.5cm]{4cm}{0cm}}
  \caption{Sample figure caption.}
  \label{fig:example}
\end{wrapfigure}
% \end{figure}


All artwork must be neat, clean, and legible. Lines should be dark enough for
purposes of reproduction. The figure number and caption always appear after the figure. The figure caption should be lowercase (except for the first word and
proper nouns); figures are numbered consecutively. Captions should be written clearly and, ideally, allow the reader to understand the key message of the figure without having to read the main text.


\paragraph{Tables.}
All tables must be centered, neat, clean, and legible.  The table number and
title always appear before the table. See~\cref{sample-table} for an example.
Captions should be written clearly and, ideally, allow the reader to understand the key message of the table without having to read the main text.
\begin{table}[h]
  \caption{Sample table title}
  \label{sample-table}
  \centering
  \begin{tabular}{lll}
    \toprule
    \multicolumn{2}{c}{Part}                   \\
    \cmidrule(r){1-2}
    Name     & Description     & Size ($\mu$m) \\
    \midrule
    Dendrite & Input terminal  & $\sim$100     \\
    Axon     & Output terminal & $\sim$10      \\
    Soma     & Cell body       & up to $10^6$  \\
    \bottomrule
  \end{tabular}
\end{table}

% % % === Introduction
\section{Introduction}
\label{sec:intro}
Introduce the problem that you are trying to solve and write 1-2 paragraphs for each of the sub-sections.  Provide a high-level summary in this section instead of detailed descriptions.
\paragraph{Problem Definition.}
While mathematical reasoning in large language models has greatly improved over recent years, it is still very prevalent in smaller, open-weight models. Models like Qwen3-4B find difficulties reasoning about quantitative accuracy, symbolic meaning, and chain-of-thought logic. This difference between models presents a clear challenge when using compact models without using any APIs or larger models to boost inference. This project works on improving the capabilities of Qwen3-4B's mathematical reasoning using methods inherent to the model.

\paragraph{Problem Significance.}
Quality mathematical inference is highly important for many LLM applications across science, math, and engineering subjects. Any logical problem-solver or data analysis application requires numerical reasoning to function. Scientific and quantitative usages are two of the most common applications for these smaller models, so it is thus important to provide high caliber mathematical reasoning for these usages.

\paragraph{Technical Challenge.}
Why is this project that you are proposing difficult? Think about challenges from (1) data processing (2) model design (3) implementation (4) engineering (5) evaluation. 

This project is challenging across multiple dimensions, from data to evaluation, because it pushes small language models to perform reliably on tasks that typically require much larger systems.

Regarding \textbf{data processing}, high-quality mathematical reasoning data is scarce and heterogeneous. Benchmarks span many different domains of math (algebra, calculus, word problems, probability) and require both free-form and multiple-choice outputs. Cleaning, normalizing, and unifying these datasets is difficult and must be done while preserving the structure needed for reasoning. Additionally, ensuring that training data is not directly related to evaluation benchmarks is nontrivial and critical for valid results.

In terms of \textbf{model design}, the main difficulty lies in enabling strong multi-step reasoning within the constraints of a small model. Unlike larger models, smaller architectures have limited capacity to implicitly learn long reasoning chains or symbolic manipulation. Designing training strategies that enhance compositional reasoning becomes a central challenge due to an inability to learn over a long period of time.

From an \textbf{implementation} standpoint, training and fine-tuning such models on diverse math tasks requires careful tokenization. Things like answer formatting and specific reasoning steps need to be exact, otherwise performance can take a serious hit. 

On the \textbf{engineering} side, careful resource management is essential. So far, training time and reproducibility are extremely relevant and important. Optimizing training pipelines and ensuring experiments are reproducible and generalizable to the whole dataset are crucial to success.

Finally, \textbf{evaluation} is particularly challenging because mathematical reasoning is sensitive not just to correctness, but to format and interpretation. Aggregating multiple benchmarks into a single unified metric requires consistent answer extraction and normalization across variable mathematic tasks. One thing we have noticed is that small models may produce either partially correct reasoning but incorrect final answers or incomplete reasoning chains. Overcoming this obstacle is vital to determining how we best measure progress, and how we can ensure fair, comprehensive, and leakage-free evaluation.


\paragraph{Contributions.}

The starter code loads up the dataset of math problems, formats them into prompts for the model, and runs a pretrained LLM on them to generate answers. It then extracts and checks those answers, scoring them and calculating accuracy. There are many limitations, but here are a few: LLM math errors, bad answer formatting, over-computation on arithmetic/logic(along with quality degradation), and repetition traps. Currently, our solution implements the following:
\begin{enumerate}
    \item \textbf{Greedy decoding} with T = 0.0 - This is standard for math reasoning, but we evaluated this experimentally to ensure it was best.
    \item \textbf{Prompt Engineering} - Added rules about keeping thinking concise(to avoid over-computation and looping), prohibiting formatting debates(to maintain a consistent format), and a decisiveness protocol(to reduce recursive self-validation).
    \item \textbf{Full set evaluation and token-truncation} - Reducing the cap on tokens allows for increased inference speed, stabilized outputs, and less formatting rambling. While we think this is detrimental for ~5\% of the problems (the most difficult ones), we think this is an overall helpful change for now.
    \item \textbf{Dataset analysis} - Failure analysis shows us that the model is doing better than it predicts. There are many false negatives presented by the evaluation pipeline, leading us to belive accuracy is higher than reported.
\end{enumerate}

Our next steps are to implement a problem classifier, to further provide prompt engineering based on problem domain(algebra, calculus, probability, etc.), and to implement advanced inference scaling, through majority voting and self-consistency. This will allow the model to explore multiple reasoning trajectories per problem and choose the majority result.

% % % === Related Work
\section{Related Work (Optional)}
\label{sec:relatedwork}
Describe related papers and contextualize your own work with respect to them. If useful, you can reference and discuss them in the introduction too. Example:
\paragraph{Related works topic 1.} \citet{radford2018improving} introduced the first GPT-1 model...  % This is an example for citing (plz remove)

\paragraph{Related works topic 2.}
% 

% % % === Methods
\section{Methods}
\label{sec:methods}
In this section, clearly describe the technical details of the problem and your approach. If needed, provide the necessary background and cite any sources that you used.


\paragraph{Problem Setting.}

The goal of this project is to improve the mathematical reasoning performance of a small (~4B parameter) language model under strict inference constraints. Given a mathematical problem, the model must generate a correct and properly formatted final answer without access to external tools or APIs. A key challenge is that the model must not only reason correctly, but also avoid common failure modes such as excessive computation, formatting errors, and truncated outputs.

\textbf{Optimization objective:}
The model is trained (or fine-tuned) to minimize the negative log-likelihood of the target output sequence:
$$\mathcal{L}_{LM}(\theta)=- \sum^N_{i=1} \sum^{T_i}_{i=1} \text{log} \space p_\theta\left( y_{i,t} | y_{i, <t}, x_i\right)$$
where $y_{i, t}$ is the $t$-th token of the sequence and $T_i$ is the sequence length.

The \textbf{dataset} is a mixture of mathematical benchmarks, spanning from high school to graduate level problems. These problems cover multiple domains of math and answer formats.

We formulate this as a supervised sequence prediction problem. We are working on a lightweight classification step for prompt selection.
\begin{enumerate}
    \item \textbf{Input:} Our input is a dataset of natural language math problems. Let a dataset consist of $N$ examples.
    $$\mathcal{D} = \{(x_i, y_i)\}^N_{i=1}$$
    The input, $x_i$, is a natural language math problem. After the classifier is completed, our input will be prompted, where $x_i$ is output from a prompt construction function, $P(x)$. $$x_i = P(x_i', c_i)$$
    where $x_i'$ is the original input and $c_i$ is the predicted problem domain(algebra, calculus, probability, etc.), which is produced by some classifier, $\mathcal{C}()$.
    $$c_i = \mathcal{C}(x_i')$$
    \item \textbf{Output:} Our output, $y_i$, is the correct answer string, typically a short mathematical expression (e.g., a number, symbolic expression, or multiple-choice label), formatted according to task requirements.
\end{enumerate}

\paragraph{Idea Summary.}
\textcolor{red}{TODO} Describe and motivate your idea to solve the problem.  How is it different from the state-of-the-art method? Provide a sketch of your idea to illustrate y our idea. Is there any risk associated with your idea? Provide alternative solutions if your idea does not work out.

% \paragraph{Method Description.}
% Describe and motivate any relevant design or architectural choices that you explored.  
% If helpful, include a picture/sketch of your approach or architecture.

% How did you split the training and validation data? How did you normalize inputs and outputs? Detail any deviations from the starter code. If valuable, you may include visualizations of the data or diagrams/sketches of your approach.



% % % === Experiments
\section{Experiments}
\label{sec:experiments}
% 
\textcolor{red}{TODO} In this section, clearly describe your empirical exploration of the problem and associated results. Provide evidence for your claims.
Analyse and discuss your results.
Note: Whenever you make claims such as \textit{Approach A performs ``better'' than approach B}, please be very specific in terms of \emph{what} it is better at and by \emph{how much}. Be transparent if it is only better on certain metrics, but not all.
\subsection{Baselines}
\label{sec:baselines}
\textcolor{red}{TODO} Describe all the baselines excluding the starter code. You should always start with simple models (such as a multi-layer Perceptron or linear regression) and gradually increase the complexity. You may use the starter code as a baseline, too.
Use an itemized list to briefly summarize each of the models, their architecture, and parameters, and provide the correct reference if possible.

We compare with the following baselines:
\begin{itemize}
    \item \textbf{Baseline A}:
    \item \textbf{Baseline B}: 
    % Etc.
\end{itemize}
% 
\subsection{Evaluation}
\label{sec:evaluation}
\textcolor{red}{TODO} Describe how you evaluate your model and baselines. There is no need to provide long explanations for the Kaggle evaluation, but clearly explain any complementary evaluations you explore in your project (if any).
% 
\subsection{Implementation details}
\label{sec:implementation_details}
\textcolor{red}{TODO} The goal of this section (together with other relevant parts of your paper) is to ensure full reproducibility without having to check your code.
Some relevant points to include:
\begin{itemize}
    \item What computational platform/GPU did you use for training and testing? How long does it take to train your model for one epoch (going through the entire training data set once)?
    \item Describe hyperparameters/regularization techniques/optimization details: What is your optimizer and learning rate, etc.? How many epochs? What is your batch size? Which HPs were tuned, and what range of values did you try? Was the tuning any different between baselines and your model? What were the final values that you used?
    \item  How did you choose them? Provide any intuition of these choices: Explain why you made these design choices. Was it motivated by your past experience? Or was it due to the limitations of your computational platform? Be transparent if any hyperparameters are inconsistent between different runs (e.g., baseline A and B were trained with different learning rates). 
\end{itemize}
% 
\subsection{Results}
\label{sec:results}
\textcolor{red}{TODO} Present the results from starter code and preliminary implementation from your idea. This can be a mix of quantitative results (e.g., validation metrics or leaderboard score) and qualitative results (e.g., example visualizations of predictions, case studies, etc.). Use figures and tables. You are welcome to include any negative findings/results.
\paragraph{Result 1.}
\paragraph{Result 2.}
% etc.
% 
% \subsection{Ablations}
% \label{sec:ablations}
% Try to explore which design choices contribute to your final performance by removing/replacing individual components. Note that ablations are most convincing when you isolate the contribution of those individual components. For example, if you are comparing architecture A and architecture B, make sure to use the same hyperparameters as much as possible (or tune them for each architecture separately and compare the best version of each). Be clear but concise.
% 
% % % === Discussion
\section{Discussion}
\label{sec:discussion}
Provide a few sentences to summarize your project. 

\paragraph{Achievements.}
\textcolor{red}{TODO} What have you achieved? What is the significance of your contribution? Did you contribute to dataset curation, algorithm improvement or model development?  Summarize your achievement in quantitative terms. We value the effort that you have devoted to this project.



% \paragraph{Lessons Learned.}
% What did you learn from this project at this point?   Brainstorm in different aspects including selecting problems, literature survey, team formulation and technical implementation. 
% If given more time, what aspects and how would you improve your project development?

\paragraph{Bottlenecks and Mitigation.}
\textcolor{red}{TODO} What are bottlenecks of the project? How did you resolve them?  Any changes that you would like to highlight from the milestone?



\paragraph{Plans Forward.}
\textcolor{red}{TODO} What do you plan to do from now until the final report? Any change of plans from the initial proposal? What decision factors contributed to this change of plans?



\paragraph{Team Responsibilities.}
\textcolor{red}{TODO} Summarize the background and expertise for each of your team member. Clearly identify each person's role and contribution in this project.

\paragraph{Timelines.}
\textcolor{red}{TODO} Estimate the timeline with the milestones to achieve your expected outcome in this class, considering the discussion provided above.
\begin{itemize}
    \item Week 2: Fill in your milestones and deliverable
    \item Week 3:
    \item ...
    \item Week 10:
\end{itemize}


\section*{LLM Usage}
The use of LLMs is allowed as a general-purpose assist tool. However,  if LLMs played a significant role in research ideation and/or writing to the extent that they could be regarded as a contributor, then authors should describe the precise role of the LLM in this section on LLM usage. 

Irrespective of the ways that LLMs were used in a given project, authors should understand that they take full responsibility for the contents written under their name, including content generated by LLMs that could be construed as plagiarism or scientific misconduct (e.g., fabrication of facts).  LLMs are not eligible for authorship.



% % % === ACKNOWLEDGMENTS AND REFERENCES START
% \section*{Acknowledgements}
% This work was supported $\dots$.
\bibliography{references.bib}
\bibliographystyle{plainnat}  % Can also use abbrvnat, unsrt, etc.
% 
% % % === APPENDIX (optional)
% \newpage
% \appendix
% \section*{Appendix}
% \section{Appendix topic 1}
% 
\end{document}